\begin{question}[index=11]
设可导函数 $y(x)$ 满足 $x^2+xy+y^2-x=0$、$y(1)=-1$，则 $\displaystyle\lim_{x\to1}\frac{(x-1)^2}{y(x)+1}=$\fillin。
\end{question}
\QuestionBreak

\begin{question}[index=12]
设 $u=f(x,y,xyz)$，函数 $z=z(x,y)$ 由方程 $\eu^{xyz}=\displaystyle\int_{xy}^z g(xy+z-t)\dd t$ 所确定，其中 $f$ 具有\allowbreak\mbox{连续偏导数}，$g$ 连续，\allowbreak\mbox{求 $x\dfrac{\partial u}{\partial x}-y\dfrac{\partial u}{\partial y}=$\fillin。}
\end{question}
\QuestionBreak

\begin{question}[index=13]
抛物面 $\Sigma\colon z=\dfrac{x^2+y^2}{2}$（$0\le z\le2$）形心的竖坐标 $\overline z=$\fillin.
\end{question}
\QuestionBreak

\begin{question}[index=14]
已知 $y_1=\eu^x-\eu^{-x}\sin2x$、$y_2=\eu^x$ 为同一二阶常系数非齐次线性微分方程的解。该方程为\fillin。
\end{question}
\QuestionBreak

\begin{question}[index=15]
设 $A$ 为 $4$ 阶方阵，$\alpha_1,\alpha_2,\alpha_3,\alpha_4$ 为线性无关的 $n$ 维向量，且
\[
\begin{aligned}
A\alpha_1&=2\alpha_1+\alpha_2+\alpha_3+\alpha_4,&
A\alpha_2&=\alpha_1+2\alpha_2+\alpha_3+\alpha_4,\\
A\alpha_3&=\alpha_1+\alpha_2+2\alpha_3+\alpha_4,&
A\alpha_4&=\alpha_1+\alpha_2+\alpha_3+2\alpha_4.
\end{aligned}
\]
则 $r(A-5E)=$\fillin.
\end{question}
\QuestionBreak

\begin{question}[index=16]
设随机变量 $X\dist N(0,\sigma^2)$，若 $X$ 落在区间 $\left(\dfrac12,1\right)$ 内的概率最大，\allowbreak\mbox{则 $\sigma^2=$\fillin。}
\end{question}
